\section{Introduction}
\label{sec:intro}

Long-horizon agents that browse, read, and reason accumulate every action--observation pair in
their context, and their accuracy degrades as trajectories grow \citep{yao2023react}.
Context folding \citep{sun2025contextfolding} lets the agent \emph{branch} into a
sub-trajectory for a localized subtask and \emph{return} with a short summary, removing the
intermediate steps from the working context. The behavior is trained end-to-end with
\foldgrpo{}, a variant of \grpo{} \citep{shao2024deepseekmath} that computes advantages over
folded contexts and adds dense process penalties. On \bcplus{}
\citep{chen2025browsecompplus}, the paper reports 0.620 pass@1 for \foldgrpo{} against 0.567
for a matched \grpo{} baseline at 36B scale.

We previously replicated the \bcplus{} track with the authors' re-implementation, substituting
\qwen{} \citep{yang2025qwen3} for \seedoss{} and a locally served \judgemodel{}
\citep{openai2025gptoss} for the proprietary judges. The mechanism replicated, the scoring
margin did not. While auditing the validation trajectories of that study we found that none of
the agent's tool observations had been tokenized correctly. The defect is specific to the
interaction between the re-implementation's turn tokenization and the \qwen{} chat template,
so it is a property of our substitution rather than of the paper's setup, but it affects every
number we had reported. Whether a bug that removes the header and the first tokens of every
observation changes the conclusions of a replication is not obvious. The trained policies had
visibly adapted to it, emptying their reasoning blocks so that observations would survive, and
an adaptation of that kind could either mask or manufacture a difference between the two
training objectives.

This paper answers that question. We characterize the bug and its mechanism, verify that the
authors' model is immune, fix the tokenization, and repeat the base-model evaluation and the
full \foldgrpo{}-versus-\grpo{} training comparison on the fixed code. The comparison uses the
same data, hyperparameters, judge, and infrastructure for both arms, with a pre-registered
temperature-1.0, four-sample evaluation tier as the decisive test. The retraining campaign was
cut short by the shared queue reclaiming the reward-serving node, so we also show how the
number of reward-bearing updates in each checkpoint can be recovered from per-step gradient
norms, and we compare the arms at matched budgets.

The contributions are: (i) a diagnosis of a template-diff tokenization bug that silently
corrupts observations for Qwen3-style templates, with a tested fix; (ii) a re-measured base
contrast showing that the fix changes behavior but not base accuracy; (iii) a repeated
\foldgrpo{}-versus-\grpo{} comparison on intact contexts, which again finds a statistical tie;
and (iv) a revision of two behavioral claims from the earlier report, namely that empty
reasoning under RL is a reward effect rather than a format artifact and that \foldgrpo{}'s
context blow-out advantage was partly a product of the bug.
